\documentclass[11pt,a4paper]{article}
\usepackage[margin=2.5cm]{geometry}
\usepackage{amsmath,amssymb}
\usepackage{booktabs}
\usepackage{array}
\usepackage{xcolor}
\usepackage{tcolorbox}
\tcbuselibrary{skins,breakable}
\usepackage{hyperref}
\usepackage[numbers,sort&compress]{natbib}
\bibliographystyle{unsrtnat}
\hypersetup{colorlinks=true,linkcolor=blue!60!black,urlcolor=blue!60!black}

\definecolor{boxblue}{RGB}{220,235,255}
\definecolor{boxorange}{RGB}{255,240,210}
\definecolor{boxgreen}{RGB}{220,245,220}

\newcommand{\term}[1]{\textbf{#1}}
\newcommand{\poly}{p}

\title{\textbf{The GTE Polynomial: A Step-by-Step Cheat Sheet}\\[6pt]
\large From Cellular Automata to the Standard Model\\[4pt]\normalsize\textit{UGP Physics Tutorial Series}
\\[2pt]\small\href{https://doi.org/\TutorialDOI}{doi.org/\TutorialDOI}}
\author{Nova Spivack\\[4pt]
\normalsize
\href{https://novaspivack.com/research}{novaspivack.com/research}\\[2pt]
\small Full monograph (P48):
\href{https://doi.org/10.5281/zenodo.20560550}{doi.org/10.5281/zenodo.20560550}
}
\newcommand{\TutorialDOI}{10.5281/zenodo.20657889}
\date{2026}

\begin{document}
\setcounter{tocdepth}{2}
\maketitle
\tableofcontents
\newpage

\begin{tcolorbox}[colback=blue!6,colframe=blue!40!black,
  title=\textbf{UGP Physics Tutorial Series}]
\textbf{Prerequisites (read first):} None — this is the first tutorial in the series.\\[4pt]
\textbf{Next tutorials:} \href{https://doi.org/10.5281/zenodo.20657919}{\emph{How the UWCA Works}} $\cdot$ \href{https://doi.org/10.5281/zenodo.20657895}{\emph{Perfect Self-Containment and MDL}} $\cdot$ \href{https://doi.org/10.5281/zenodo.20657876}{\emph{How MDL Selects the 19-Bit Polynomial}}\\[4pt]
\textbf{See also:} \href{https://doi.org/10.5281/zenodo.20657913}{\emph{From the Arithmetic Substrate to Particle Masses}}
\end{tcolorbox}


%─────────────────────────────────────────────────────────────────────────────
\section{The Basics: What Kind of Object Is This?}
%─────────────────────────────────────────────────────────────────────────────

\begin{tcolorbox}[colback=boxblue,colframe=blue!40!black,title=The one-line answer]
$p$ is a cellular automaton rule that tells you the new state of a cell from
its two neighbours and its own current state.  All arithmetic is done on a
clock with 7 positions instead of 10.
\end{tcolorbox}

\subsection{Cellular automata in one sentence}

A cellular automaton is a row of cells, each holding a value.  At every
time step, \emph{every} cell simultaneously updates its value using a fixed
rule that looks only at the cell itself and its nearest neighbours.
Stephen Wolfram's Rule 110 is a well-known example with 2 possible cell
states (0 or 1).

\subsection{The GTE polynomial has 7 states}

In the GTE framework, each cell holds one of \textbf{7} possible states:
$\{0, 1, 2, 3, 4, 5, 6\}$.  This makes it a \emph{k\,=\,7}, \emph{r\,=\,1}
rule in Wolfram's notation:
\begin{itemize}
  \item $k = 7$: seven possible states per cell.
  \item $r = 1$: the rule looks one neighbour to each side (Left and Right).
\end{itemize}

\subsection{The truth table}

Because each of the three inputs (Left, Center, Right) can be any of
7 values, there are $7^3 = 343$ possible input combinations.  A complete
specification of the rule is a table with 343 rows, one output per row.

The GTE polynomial is a compact algebraic formula that reproduces all
343 rows of this truth table in a single expression.

%─────────────────────────────────────────────────────────────────────────────
\section{The Arithmetic: What Is GF(7)?}
%─────────────────────────────────────────────────────────────────────────────

\begin{tcolorbox}[colback=boxorange,colframe=orange!60!black,title=Key idea]
GF(7) is \emph{clock arithmetic} with 7 positions instead of 12.
After every calculation you wrap around: if the result reaches 7, it
becomes 0; if it goes to 8, it becomes 1; and so on.
\end{tcolorbox}

\subsection{Clock arithmetic (modular arithmetic)}

On a 12-hour clock, $11 + 3 = 2$ (not 14).  In GF(7), the same idea applies
but with 7 positions.  Formally: take any integer result and compute its
\term{remainder when divided by 7}.

\bigskip
\begin{center}
\begin{tabular}{lll}
\toprule
Ordinary result & Divided by 7 & Remainder (mod 7) \\
\midrule
$6$  & $7 \times 0 + 6$ & $6$ \\
$7$  & $7 \times 1 + 0$ & $0$ \\
$8$  & $7 \times 1 + 1$ & $1$ \\
$15$ & $7 \times 2 + 1$ & $1$ \\
$-1$ & $7 \times (-1) + 6$ & $6$ \\
\bottomrule
\end{tabular}
\end{center}
\bigskip

So $-1$ and $6$ are the \emph{same element} in GF(7) — they differ by a
multiple of 7, which disappears under mod 7.

\subsection{Why it is called a ``field''}

A \term{field} is any collection of numbers where you can add, subtract,
multiply, and divide (except by zero) and the usual rules of algebra hold
(commutativity, associativity, distributivity).  The rational numbers are a
field; the real numbers are a field; GF(7) is a field.

The crucial property: every nonzero element in GF(7) has a
\term{multiplicative inverse} — a number you can multiply it by to get~1.

\begin{center}
\begin{tabular}{ccc}
\toprule
Element $x$ & Its inverse $x^{-1}$ & Check: $x \times x^{-1} \bmod 7$ \\
\midrule
1 & 1 & $1\times1=1$ \\
2 & 4 & $2\times4=8\equiv1$ \\
3 & 5 & $3\times5=15\equiv1$ \\
4 & 2 & $4\times2=8\equiv1$ \\
5 & 3 & $5\times3=15\equiv1$ \\
6 & 6 & $6\times6=36\equiv1$ \\
\bottomrule
\end{tabular}
\end{center}

This means you can always ``divide'' in GF(7) by multiplying by the
inverse — so all of standard polynomial algebra carries over.

The name \emph{Galois Field} honours Évariste Galois, the 19th-century
mathematician who proved that such finite fields exist for every prime
(7 is prime) and for prime powers, and for nothing else.

%─────────────────────────────────────────────────────────────────────────────
\section{The Polynomial}
%─────────────────────────────────────────────────────────────────────────────

\begin{tcolorbox}[colback=boxgreen,colframe=green!50!black,
  title=The GTE master polynomial]
\[
  p(L,\,C,\,R) \;=\; C + R - CR - LCR \pmod{7}
\]
\begin{center}
$L$ = Left neighbour,\quad $C$ = Center cell,\quad $R$ = Right neighbour\\
All values in $\{0,1,2,3,4,5,6\}$;\quad all arithmetic mod 7.
\end{center}
\end{tcolorbox}

\subsection{Reading the expression term by term}

Written out in plain English, the output equals:
\[
  \underbrace{C}_{\text{center value}}
  +\;\underbrace{R}_{\text{right value}}
  -\;\underbrace{C \times R}_{\text{centre--right product}}
  -\;\underbrace{L \times C \times R}_{\text{three-body product}}
\]
then reduce mod 7.

\bigskip
\begin{center}
\begin{tabular}{clll}
\toprule
Term & Plain English & Coefficient & Effect \\
\midrule
$C$   & center cell value & $+1$ & carries the center forward \\
$R$   & right cell value  & $+1$ & adds the right neighbour \\
$CR$  & center $\times$ right & $-1\equiv 6$ & pairwise interaction, subtracted \\
$LCR$ & left $\times$ center $\times$ right & $-1\equiv 6$ & three-body interaction, subtracted \\
\midrule
$L$   & left cell value   & $0$  & absent: left alone has no direct effect \\
$LC$  & left $\times$ center & $0$ & absent \\
$LR$  & left $\times$ right  & $0$ & absent \\
\bottomrule
\end{tabular}
\end{center}

\bigskip
\noindent
The notation $-CR$ means ``subtract $C$ times $R$''; equivalently, add
$6 \times CR$ (since $-1 \equiv 6$ in GF(7)).  Both forms are correct; the
signed form is the conventional mathematical shorthand.

\subsection{Why ``multilinear'' and ``degree 3''}

The polynomial is \term{multilinear}: each input variable ($L$, $C$, $R$)
appears at most to the first power — no $C^2$ or $L^3$ terms.  This is a
special property of this particular rule.  It means the polynomial is
completely determined by its values at the $2^3 = 8$ binary input points
$\{0,1\}^3$ — which is the key to the uniqueness theorem.

The \term{degree} of the polynomial is 3 (the highest-degree term is $LCR$,
which has degree $1+1+1=3$).  A degree-$\leq 3$ multilinear polynomial
over three variables has at most $2^3 = 8$ independent coefficients
(one per subset of $\{L,C,R\}$, including the empty set).  The GTE polynomial
uses 4 of those 8 slots; the other 4 have coefficient 0.

\subsection{Computing a value: worked example}

\textbf{Inputs:} $L = 1$, $C = 3$, $R = 2$.

\[
  p(1,3,2) = 3 + 2 - (3\times2) - (1\times3\times2)
           = 3 + 2 - 6 - 6
           = -7
           \equiv 0 \pmod{7}
\]

The new center-cell value is \textbf{0}.

\bigskip
\textbf{Inputs:} $L = 0$, $C = 1$, $R = 1$.

\[
  p(0,1,1) = 1 + 1 - (1\times1) - (0\times1\times1)
           = 1 + 1 - 1 - 0
           = 1
\]

The new center-cell value is \textbf{1}.

%─────────────────────────────────────────────────────────────────────────────
\section{The Connection to Rule 110}
%─────────────────────────────────────────────────────────────────────────────

\subsection{Restricting to binary inputs}

If you only ever feed the polynomial inputs from $\{0,1\}$ — ignoring states
2 through 6 — then the output is also always in $\{0,1\}$ (you can verify
this by checking all 8 cases).  This binary restriction is a standard
2-state, 3-neighbour cellular automaton rule.

\begin{center}
\begin{tabular}{ccc|cc}
\toprule
$L$ & $C$ & $R$ & $p(L,C,R)$ & Rule 110? \\
\midrule
0 & 0 & 0 & 0 & 0 \\
0 & 0 & 1 & 1 & 1 \\
0 & 1 & 0 & 1 & 1 \\
0 & 1 & 1 & 1 & 1 \\
1 & 0 & 0 & 0 & 0 \\
1 & 0 & 1 & 1 & 1 \\
1 & 1 & 0 & 1 & 1 \\
1 & 1 & 1 & 0 & 0 \\
\bottomrule
\end{tabular}
\end{center}

Reading the output column from bottom to top gives the binary number
$01101110_2 = 110_{10}$ — the Wolfram number of \textbf{Rule 110}.  The
binary restriction of $p$ is Rule 110 exactly.  This is a theorem, not a
coincidence.

\subsection{Why Rule 110 matters}

Rule 110 is the simplest known cellular automaton that is
\term{computationally universal}: given the right initial pattern, it can
simulate any Turing machine.  Matthew Cook proved this in 1994 (published
2004).  Wolfram championed it as evidence that simple rules can produce
arbitrarily complex behaviour.

Because $p$ restricts to Rule 110, the GTE framework~\cite{SpivackGTECompleteFramework} inherits that
universality: the GTE substrate can simulate any computation.

\subsection{Why GF(7) and not just Rule 110?}

Rule 110 alone costs 8 bits and certifies exactly one thing: Turing universality.
It cannot certify particle generations, charge quantization, color symmetry,
chirality, or baryon number, because a 2-state binary field has no structure to
express them.
The GTE polynomial is a $\mathrm{GF}(7)$ rule whose binary restriction happens to
be Rule 110 --- Rule 110 is the \emph{shadow} of $p$ on binary inputs, not its
source.
The additional 11 bits (over the 8-bit Rule 110 floor) encode the $\mathrm{GF}(7)$
field structure that makes five further physical certifications possible:
\begin{enumerate}
  \item \textbf{Generation arithmetic.} The GTE orbit $\{73, 42, 275\}$ requires
    7 distinct states to track distinct winding numbers; a binary field collapses
    all three orbit types to a single class.
  \item \textbf{Fractional electric charge.} Charges $\tfrac{1}{3}$ and
    $\tfrac{2}{3}$ arise from $\mathbb{Z}_7$ winding ratios in
    $\mathrm{GF}(7)^*$; binary fields produce only integer charges mod 2.
  \item \textbf{Colour charge ($\mathbb{Z}_3$).} The Sylow-3 subgroup of
    $\mathrm{GF}(7)^* \cong \mathbb{Z}_6$ provides the $\mathrm{SU}(3)$ colour
    symmetry.  No $\mathbb{Z}_3$ subgroup exists in any binary field.
  \item \textbf{Chirality (V$-$A).} The chiral $\mathbb{Z}_2$ subgroup
    $\{1, 6\} \subset \mathrm{GF}(7)^*$ distinguishes forward from mirror orbits.
    Binary systems are symmetric under all reflections.
  \item \textbf{Baryon number.} Baryon number is conserved from a $\mathbb{Z}_7$
    winding conservation law on the GTE orbit; it is not definable in a 2-state
    field.
\end{enumerate}
More compact does not mean more powerful: a 19-bit $\mathrm{GF}(7)$ polynomial
does strictly more physics than any 8-bit binary rule.

%─────────────────────────────────────────────────────────────────────────────
\section{The Uniqueness Theorem}
%─────────────────────────────────────────────────────────────────────────────

\subsection{The question}

There are $7^{343} \approx 10^{290}$ possible k=7, r=1 cellular automaton
rules.  Why is $p$ the right one?

\subsection{Two independent selection paths}

\textbf{Path A (spectrum arm): how the orbit pins the coefficients.}

\medskip
\noindent\textit{Step 1 — The GTE orbit.}
The UGP arithmetic starts from the Lepton Seed $(1, 73, 823)$ and applies
the map $T$ twice, producing the three $N_{\mathrm{eff}}$ values
$\{73, 42, 275\}$.  These together with intermediate triples form the
\term{GTE orbit}: a specific sequence of integer triples over $\mathbb{Z}$.

\medskip
\noindent\textit{Step 2 — The parity shadow.}
Reduce every coordinate of every orbit triple mod 2 — keep only whether each
value is even (0) or odd (1).  This produces a sequence of binary triples.
Collecting the unique ones gives 7 distinct binary $(L, C, R)$ patterns,
each paired with an output bit also read from the orbit.  These are 7 rows
of a binary truth table, filled in by the orbit's own arithmetic — no choice
involved.

\medskip
\noindent\textit{Step 3 — The 8th row.}
One binary input pattern is not covered by the parity shadow: $(1, 1, 1)$
(all three neighbours odd).  This is filled in by the
\term{vacuum-transparency condition}: the all-zeros state $(0,0,0)$ must be
a fixed point of the dynamics (the vacuum stays vacuum).  Evaluating $p(0,0,0)=0$
and propagating this constraint forces the $(1,1,1)$ output to be 0.
Now all 8 binary rows are specified.

\medskip
\noindent\textit{Step 4 — Polynomial interpolation.}
A multilinear polynomial of degree $\leq 3$ over three binary variables has
exactly $2^3 = 8$ basis monomials:
\[
  1,\; L,\; C,\; R,\; LC,\; LR,\; CR,\; LCR.
\]
With 8 known input--output pairs and 8 unknown coefficients, this is a
standard linear system $\mathbf{V}\,\mathbf{c} = \mathbf{y}$, where
$\mathbf{V}$ is the $8\times8$ matrix of monomial evaluations at the binary
points.  This matrix is a Vandermonde-type matrix that is invertible over
$\mathrm{GF}(7)$ (confirmed by the theorem).  Solving the system over
$\mathrm{GF}(7)$ gives a unique coefficient vector $\mathbf{c}$.

\medskip
\noindent\textit{Step 5 — The result.}
The unique solution is:
\[
  (c_1,\, c_L,\, c_C,\, c_R,\, c_{LC},\, c_{LR},\, c_{CR},\, c_{LCR})
  \;=\; (0,\; 0,\; 1,\; 1,\; 0,\; 0,\; 6,\; 6)
\]
which — remembering that $6 \equiv -1$ in $\mathrm{GF}(7)$ — gives:
\[
  p(L,C,R) \;=\; C + R - CR - LCR \pmod{7}.
\]
The polynomial is not chosen; it is the \emph{unique} solution to the linear
system imposed by the orbit's parity shadow.  Different orbit $\Rightarrow$
different right-hand side $\Rightarrow$ different (or no valid) polynomial.
The GTE orbit happens to produce exactly the Rule-110 truth table on the
right-hand side — and that is the content of the theorem.

\textbf{Path B (MDL arm).}
Among all k=7, r=1 rules consistent with the physical constraints (the
Standard Model derivability criterion), which has the shortest description?
Specifying $p$ as a polynomial requires only \textbf{19 bits}:
\begin{itemize}
  \item 8 bits — the Rule 110 truth table (the binary restriction)
  \item 3 bits — the prime modulus ($\lceil\log_2 7\rceil = 3$)
  \item 5 bits — the monomial selector (which of the 8 basis terms are
    nonzero)
  \item 1 bit — chirality (left--right orientation)
  \item 2 bits — the $-1$ coefficient signs
\end{itemize}
A naive lookup table for the same rule requires $343 \times \log_2 7 \approx
963$ bits — about 50 times larger.  The MDL uniqueness theorem (machine-certified
in Lean 4, zero sorry, zero custom axioms) proves no shorter valid
specification exists.

\textbf{The closing edge.}
Paths A and B converge on the same polynomial.  The
\term{Direct-Interpolation Lift Theorem} (machine-certified) proves that the
orbit's parity shadow, combined with MDL minimality, uniquely interpolates
$p$ from among all $10^{290}$ candidates.  The circuit is closed.

%─────────────────────────────────────────────────────────────────────────────
\section{What the Polynomial Produces}
%─────────────────────────────────────────────────────────────────────────────

Starting from $p$ as the algebraic certificate, with zero free parameters and
no fitting to experimental data, the GTE framework derives:

\begin{itemize}
  \item All 9 charged-fermion masses (0.295\% RMS vs PDG 2022~\cite{SpivackSM})
  \item The electroweak mixing angle (Weinberg angle)
  \item The strong CP angle $\theta_{QCD} = 0$ exactly (three independent proofs)
  \item The fine-structure constant $\alpha^{-1} = 2^7 + 3^2 = 137$ at tree level
  \item 3+1 spacetime dimensions (forced, machine-certified)
  \item Einstein's equations from a variational principle
  \item Newton's constant $G_N$ to 0.040\% from orbit arithmetic
  \item The Born rule (four independent derivations, two machine-certified)
  \item The cosmological constant bracket (enclosing the Planck 2018 value)
  \item $40+$ further observables across all Standard Model sectors
\end{itemize}

The derivation chain: $p$ is the algebraic certificate; it certifies what the
continuum physical field $\Phi_{\rm MDL}$ (the actual universe) must contain.
The Algebraic Lifting Theorem bridges the two: every algebraic fact proved in the
certificate holds in the physical field.  Physical predictions --- masses, forces,
constants --- live in the field; the polynomial certifies them.
The full derivation machinery is documented in 52 papers, with over 300 Lean~4
modules (all zero sorry, zero custom axioms), all publicly available at
\url{https://zenodo.org/records/20560550}.

\bigskip
\begin{tcolorbox}[colback=boxorange,colframe=orange!60!black,
  title=Special numbers: cyclotomic identities]
The GTE Higgs coefficient $c_H = 13$ satisfies a triple identity:
\[
  c_H \;=\; \Phi_{12}(2) \;=\; F_7 \;=\; 13,
\]
where $\Phi_{12}(n) = n^4 - n^2 + 1$ is the twelfth cyclotomic polynomial
and $F_7 = 13$ is the 7th Fibonacci number
$(F_1=1,\,F_2=1,\,F_3=2,\,F_4=3,\,F_5=5,\,F_6=8,\,F_7=13)$.

The same cyclotomic polynomial at $n = 3$ gives the Lepton Seed entry:
$\Phi_{12}(3) = 3^4 - 3^2 + 1 = 81 - 9 + 1 = 73$.

As a consequence, the tau-lepton $N$-value is
$b_3 = b_2 + F_{c_H} = 42 + F_{13} = 42 + 233 = 275$,
recovering the GTE orbit triple $\{73, 42, 275\}$.
Furthermore, $F_{13} \bmod 7 = 233 \bmod 7 = 2$, the up-quark winding class.

The appearance of $F_7 = 13$ as the quotient gap is not merely a numerical
coincidence: $|q_2 - q_1| = F_{\mathrm{ord}_7(2) + N_c + 1} = F_7$ is the
\emph{quotient gap theorem}, proved from the $\mathbb{Z}_7 \times \mathbb{Z}_3$
structure (machine-certified in Lean~4, zero sorry, zero custom axioms,
\texttt{gt012\_gap\_from\_z7\_z3}).
Fibonacci numbers enter the GTE arithmetic because the algebraic structure of
the theory forces them.
\end{tcolorbox}

\bigskip
\begin{tcolorbox}[colback=boxblue,colframe=blue!40!black,
  title=One-paragraph summary]
\textbf{$p(L,C,R) = C + R - CR - LCR$} is a 7-state, radius-1 cellular
automaton update rule.  All arithmetic is mod 7 (clock arithmetic with 7
positions).  Its inputs are three integers from 0 to 6 (the cell's left
neighbour, itself, and its right neighbour); its output is an integer from 0
to 6 (the new center value).  Restricted to binary inputs, it is Rule 110
exactly.  It is selected by minimum description length from $10^{290}$
candidates, and this selection is forced — the Standard Model's own algebraic
structure, when asked ``what is the shortest complete description of the GTE
orbit?'', hands back this polynomial.  It is the algebraic certificate for
the continuum field $\Phi_{\rm MDL}$ (the physical universe), certifying via
the Algebraic Lifting Theorem that $\Phi_{\rm MDL}$ must contain the masses,
forces, geometry, and measurement structure of the Standard Model --- derived,
not fitted.
\end{tcolorbox}


\section*{Further Reading}
\begin{tcolorbox}[colback=orange!8,colframe=orange!50!black,
  title=\textbf{Primary Sources}]
The results in this tutorial are derived in the following papers,
all available open-access on Zenodo and at
\href{https://novaspivack.com/research}{novaspivack.com/research}:
\begin{itemize}
  \item \textbf{P48 --- The Complete GTE Framework} (main synthesis monograph):\\
        \href{https://doi.org/10.5281/zenodo.20560550}{doi.org/10.5281/zenodo.20560550}
  \item \textbf{P01 --- Standard Model Parameters from the UGP} (first-principles derivation):\\
        \href{https://doi.org/10.5281/zenodo.20168787}{doi.org/10.5281/zenodo.20168787}
  \item \textbf{P40 --- GF(7) Polynomial Universality} (uniqueness theorem):\\
        \href{https://doi.org/10.5281/zenodo.20417566}{doi.org/10.5281/zenodo.20417566}
\end{itemize}
Lean~4 machine proofs: \href{https://github.com/novaspivack/ugp-lean}{github.com/novaspivack/ugp-lean}
\end{tcolorbox}

\bibliography{../../papers/bib/Spivack_Papers_Bibliography}
\end{document}
